\section{Related Work}

Our study sits at the intersection of four research areas: quality-diversity and diversity in procedural content generation, stochasticity and prompt sensitivity in language models, reproducibility in generative systems, and design space theory. We situate same-input divergence within each area and clarify the experimental and conceptual distinctions that make our contribution novel.

\subsection{Quality-Diversity and Diversity in Procedural Generation}

A significant strand of computational creativity and procedural content generation (PCG) research treats diversity as a primary design objective. The MAP-Elites algorithm of \citet{mouret2015illuminating} exemplifies this orientation: it maintains an archive of solutions partitioned by behavioral descriptors, explicitly maximizing both quality within each niche and variety across niches. This quality-diversity (QD) framing has been widely adopted in game content generation \cite{gravina2019quality}, where diverse outputs are desirable to avoid player boredom and cover the space of possible experiences.

Diversity measurement in PCG follows a complementary research agenda. \citet{smith2010analyzing} propose expressive range analysis, a framework for visualizing and quantifying the variety of outputs a generator can produce across dimensions such as linearity and leniency. \citet{shaker2014evolving} apply evolutionary search to generator parameters with diversity as a fitness component, treating expressiveness as a property to be maximized through search.

Our work differs from this tradition in a fundamental respect. The QD and expressive range literatures ask: \textit{how much diversity can this generator produce, and how do we maximize or measure it?} We ask a different question: \textit{given a fixed input specification, which components of the output are diverse and which are not, and why?} We do not seek to maximize divergence; we seek to characterize its structure. The distinction is between diversity as a goal to engineer and divergence as a phenomenon to explain. Understanding this structure has immediate consequences for how practitioners deploy AI creative generators in production pipelines, because it determines where human oversight is necessary and where a single AI run can be trusted.

\subsection{Stochasticity and Prompt Sensitivity in Language Models}

Language models introduce stochasticity through temperature-controlled sampling, and several lines of research examine how this stochasticity shapes output variation. \citet{ficler2017controlling} study the effect of stylistic parameters, including temperature, on the diversity and quality of generated text, establishing that temperature is a primary lever for controlling output variation. This work treats temperature as a design parameter to be set; our study treats temperature as fixed and asks what remains variable despite that setting.

A related and more recent body of work examines prompt sensitivity: the observation that LLM outputs vary substantially with small changes to the input prompt. \citet{jiang2020can} demonstrate that different phrasings of the same factual query yield dramatically different model responses, coining the term ``prompting'' as a search problem. \citet{zhao2021calibrate} show that LLM classification is sensitive to label ordering and demonstration selection in few-shot settings, and \citet{lu2022fantastically} extend this to show that permutation of in-context examples strongly affects output quality, with optimal orderings differing across models and tasks.

These prompt sensitivity studies vary the input to measure how sensitive the output is to input perturbation. Our experimental paradigm is precisely the inverse: we \textit{fix the prompt} and measure output variation across independent runs. This targets a distinct phenomenon --- the natural stochasticity inherent in a generative model at fixed temperature, rather than sensitivity to prompt design. Moreover, our analysis operates at the level of a structured, multi-stage creative artifact (a puzzle hunt with rounds, answer words, mechanisms, and clue formulations), not at the level of a single model response. This multi-level analysis is necessary to ask which \textit{components} of the output are sensitive to stochasticity, a question that single-response studies cannot address.

\subsection{Reproducibility in Generative Systems}

A third relevant area is reproducibility in machine learning and generative systems. \citet{pineau2021improving} provide a comprehensive framework for reproducibility in ML research, emphasizing seed control, hyperparameter reporting, and environment specification as prerequisites for replicable results. The general thrust of this literature is that stochasticity is a confound to be controlled: reproducibility is achieved when two runs of the same system produce the same output.

In computational creativity, \citet{colton2011computational} examine reproducibility as one criterion for evaluating creative systems, raising the question of whether a system's outputs are predictably related to its inputs. In PCG specifically, seed control is standard practice: fixing the random seed guarantees that a procedural generator produces identical output on repeated runs, which is essential for debugging, level handcrafting, and player-facing consistency.

Our study takes the opposite stance toward stochasticity. Rather than eliminating or controlling it, we deliberately allow it to operate and use the resulting variation as data. Our goal is not to make runs identical but to understand \textit{which decisions are inherently stable across runs and which are inherently variable}. Stable decisions are those where the design space is sufficiently constrained that independent runs converge to the same choice; variable decisions are those where the design space offers many equally valid options and independent runs do not converge. This framing converts stochasticity from a nuisance to be suppressed into an instrument for characterizing the structure of the design space itself.

\subsection{Design Space Theory and Constrained Creative Decisions}

The stable/variable distinction we observe empirically connects to a principled theoretical tradition in design cognition. \citet{simon1969sciences} characterizes design as search through a space of possibilities under constraints: the designer's task is to find an acceptable point in a large space of potential solutions, guided by problem-structure and satisfaction criteria. The size and structure of the relevant design space determines how constrained any particular decision is.

\citet{lawson1997designers} extends this framework with empirical observations of design cognition, distinguishing between problem-focused strategies (working from constraints inward to solutions) and solution-focused strategies (proposing solutions and testing them against constraints). Critically, this work shows that experienced designers quickly identify which aspects of a design problem are heavily constrained (leaving little creative latitude) and which are relatively free (requiring deliberate creative choice). The stable decisions we observe --- round structure, puzzle type variety, meta mechanism type --- are analogous to heavily constrained decisions in Lawson's framework: domain knowledge, audience requirements, and structural specifications together reduce the space of valid options to a small set, and independent agents navigating that space converge on similar choices. The variable decisions --- specific answer words, extraction mechanism details within a puzzle type, individual clue formulations --- correspond to design-space-free decisions: many options are equally valid, no constraint privileges one over another, and independent agents diverge.

\citet{gero1990design} formalizes this distinction in the function-behavior-structure model of design, which separates the functional requirements a design must satisfy (analogous to our domain constraints), the behavioral properties it exhibits (observable outputs), and the structural decisions that produce those behaviors (analogous to our design-space choices). Structural decisions that are over-determined by functional requirements leave little latitude and tend to be stable; structural decisions with many behavioral-equivalent options are inherently variable. Our empirical findings in a large-scale AI creative pipeline provide concrete, quantitative support for this theoretical prediction.

\subsection{Same-Input Divergence as a Computational Creativity Criterion}

Characterising the reliability structure of a creative pipeline --- identifying which generative decisions converge across independent runs and which diverge --- is a prerequisite for evaluating that pipeline as a creative system under established CC frameworks. The SPECS framework of \citet{jordanous2012standardised} provides a standardised procedure for CC evaluation, requiring that evaluation criteria be applied consistently and that a system's outputs be assessable as creative artefacts. This evaluation is only possible if we can distinguish systematic output properties from stochastic noise: a system whose every run produces a categorically different meta mechanism (Level~1) and zero shared answer words (Level~3) cannot be assessed under criterion C8 (``the system reliably produces creative artefacts'') without first characterising which decisions drive that instability. By contrast, a system where panel scores are stable across runs (Level~6) and pass/fail verdicts converge satisfies C8 in the relevant dimension. Same-input divergence operationalizes this criterion: it makes C8 empirically testable by decomposing pipeline reliability into per-decision measurements.

This framing connects to \citet{ritchie2007some}'s evaluative criteria for creative systems, which require that artefact quality be reproducible across independent runs for quality assessments to be meaningful. Ritchie's criteria cannot be applied if the relationship between input specification and output quality distribution is unknown. The same-input divergence paradigm provides the measurement instrument needed to establish this relationship: by running the pipeline twice from identical inputs and comparing outputs at multiple abstraction levels, we identify which quality dimensions are stable (and therefore assessable as reliable system properties) and which are contingent on incidental generative choices. We make no claim that divergence is evidence of creativity, nor that convergence implies its absence. Rather, we argue that characterising the structure of same-input divergence is a prerequisite for any CC evaluation of a multi-stage generative pipeline --- you cannot evaluate creative reliability without first measuring it.

\subsection*{Gap Statement}

Taken together, these four research areas provide partial predecessors but no direct treatment of the problem we address. Quality-diversity work maximizes or measures diversity as a goal but does not ask where divergence originates in a multi-stage pipeline or why some stages converge while others do not. Prompt sensitivity work varies inputs to measure output sensitivity but does not fix the input and characterize natural stochasticity across a structured creative artifact. Reproducibility work suppresses stochasticity rather than characterizing its structure. Design space theory provides the conceptual vocabulary for constrained vs.\ unconstrained decisions but has not been connected empirically to the behavior of LLM-based creative generators.

No prior work characterizes the structure of same-input divergence at multiple levels of abstraction in an AI creative pipeline. Existing work either maximizes diversity, controls it, or measures it holistically without decomposing where in a multi-stage generative process it originates and which stages are responsible. Our study fills this gap with an empirical same-input experiment, a five-level comparison methodology, and a principled explanation of the resulting divergence structure grounded in design space theory.
